\documentclass[12pt]{article}
\usepackage[utf8]{inputenc}
\usepackage[T1]{fontenc}
\usepackage{hyperref}
\title{Chin, Rawngtu Language Resources for K-12 Curriculum Development}
\author{Compiled by Claude Sonnet 5.5 (claude-sonnet-5-5) for Andr\'as Kornai}
\date{}
\begin{document}
\maketitle

\section{Introduction}
Rawngtu Chin (ISO 639-3 \texttt{weu}; also called Welaung) is a Kuki-Chin language of western Myanmar. It is spoken mainly in about ten villages, most within Matupi and Mindat townships in Chin State, with a smaller group in Magway Region. Estimates range from about 5,000 speakers (SIL, 2008) to a Joshua Project population figure of 5,600. The family given in the task prompt (``Bookkeeping'') is evidently an error. The sources I found place the language in Sino-Tibetan (Kuki-Chin). Wikipedia calls it a ``purported'' language, and Glottolog is said to flag it as possibly not distinct, while Joshua Project and Asia Harvest treat it as a distinct people group.

\section{Primary Educational Resources}
I found no primers or school materials. Asia Harvest reports that the language had no written form until a Baptist pastor created a script around the start of this century.

\section{Linguistic and Reference Materials}
\begin{itemize}
\item Welaung language (Wikipedia): \url{https://en.wikipedia.org/wiki/Welaung_language}
\item Joshua Project profile (four dialects, about 90\% lexical similarity among them; lexical similarity of 84\% with Matu and 67\% with Kaang): \url{https://joshuaproject.net/people_groups/15902}
\end{itemize}

\section{Religious and Community Resources}
Most speakers are Christian, and Scripture portions were reportedly published in 2001. Asia Harvest profile: \url{https://www.operation.asiaharvest.org/MYANMAR-Rawngtu.pdf}

\section{Digital and Audio Resources}
I found none specific to Rawngtu.

\section{Curriculum Development Context}
The language is stable, but little is documented and it has no known school use. Many speakers also use Burmese and Muun. Any curriculum would need community-led orthography work, perhaps building on the script created by the Baptist pastor and the 2001 Scripture portions.
\end{document}